\label{part:physical-isomorphism}

我们在前两篇建立了结构、内容及其联合状态的表达。本篇进一步说明这些内部变量怎样与任务环境连接。对应从观测开始，经由对象估计和认知组织，进入预测、推理与行动；每一步都有自己的条件与误差。我们不要求内部模型复制环境的全部状态，而要求它保留当前任务所需的关系，并在新证据出现时能够修订。

本篇先讨论测量、对象身份、粗粒化与领域约束，再讨论语言、视觉和推理如何形成可维护的认知对象。随后，我们考察不同领域在变量类型、时间尺度与有效性标准上的差异，并建立跨域转换与调度的接口。最后，我们用明确的映射说明状态压缩何时能够保持转移，推导近似误差怎样传播，区分任务有效、行为等价与严格结构对应。

这些讨论使智能的一般理论与具体领域知识保持清楚关系。环境中的机械规律、实验机制或社会规则进入相应模型；内部认知活动则由状态组织、证据更新与任务接口描述。HSF-HD 提供一般的分析语言，AgentOS 在其下实现对象记录、工具协作和运行流程。我们以可检查的条件连接两者，而不把形式相似性当成全部机制已经相同的证明。
